\section{Par de lectores genealógicos sobre el atlas completo}
\label{sec:operadores-masa-two-way-rev11}
\index{masa!operadores bidireccionales}
\index{atlas!lectores cronológico y dodecafásico}

La ruta enriquecida posee dos representaciones internas complementarias,
\[
 \Gamma\longmapsto
 \bigl(\gamma_{108}(\Gamma),\tau_{12}(\Gamma)\bigr)
 \in\mathbb F_3^{108}\times\{-1,0,+1\}^{12}.
\]
La primera conserva la historia trítica completa; la segunda conserva el
registro dodecafásico firmado. Cada proyección induce un lector natural y ambos
se evalúan antes de abrir masas, etiquetas PDG o residuales.

\HMTClaim{MD-R-MASS-TWOWAY-READER-PAIR}
\subsection{Par de lectores genealógicos \(\mathbf K(\Gamma)=(K_D(\Gamma),K_\Omega(\Gamma))\): dominio, codominio y evaluaciones escalares}

Con la notación de las
\cref{eq:KD-rev11,eq:kappaD-rev11,eq:kappa-omega-rev11}, definimos
\begin{equation}
 \mathbf K(\Gamma)
 :=\bigl(K_D(\Gamma),K_\Omega(\Gamma)\bigr)
 \in\left(\mathbb Z^2\times\tfrac12\mathbb Z\right)
 \times\mathbb Z^3,
 \label{eq:lector-conjunto-masas-rev11}
\end{equation}
donde
\[
 K_D=\left(D_{27}+D_{36},D_1,\frac{D_4-D_3}{2}\right),
 \qquad
 K_\Omega=(\Omega_{90\mid120},54,P_6).
\]
Las evaluaciones escalares son
\begin{align}
 \kappa_D(\Gamma)
 &=D_{27}+D_{36}-\alpha_{\rm HMT}D_1
 +\frac{\Delta_4}{2}(D_4-D_3),
 \label{eq:kappaD-general-rev11}\\
\kappa_\Omega(\Gamma)
&=\Omega_{90\mid120}(\tau_\Gamma)
-54\alpha_{\rm HMT}+P_6(\tau_\Gamma)\Delta_4.
\label{eq:kappaOmega-general-rev11}
\end{align}
La segunda coordenada exige una distinción causal. En APP, el salto
espectral primordial es \(9-5=4\); el coeficiente local \(2\), la compresión
\(6=51-45\) y el despliegue bidimensional \(54=9\cdot6=459-405\) son
invariantes posteriores y tipados, conforme al
\cref{thm:app-jerarquia-cuatro-dos-seis-cincuentaycuatro}. En el lector
cronológico de desplazamientos, en cambio, la coordenada electrónica es
\(B_D=D_1(\Gamma_e)=54\): cuenta fronteras inmediatas y coincide con la
periodicidad cinemática de orden \(54\) del endomorfismo CT108. El
\(54\) fijo de
\(K_\Omega\) es la coordenada común transferida al hacer compatibles ambos
lectores en la reducción electrónica. Su igualdad numérica con el defecto
global de APP constituye un control cruzado posterior; no convierte al
número \(54\) en la diferencia primordial entre suma y producto.

La igualdad \(4\cdot90=3\cdot120=360\) explica la genealogía común de
los dos lectores. \(K_D\) compara desplazamientos de la palabra trítica
\(\gamma_{108}\);
\(K_\Omega\) compara energías de ventanas en la palabra de doce fases. La
coincidencia de origen no identifica sus cocientes de información.
En el dominio certificado \(\Gamma_{324}^{\rm or}\), la periodicidad de
media vuelta hace pares todos los \(D_k\), de modo que
\((D_4-D_3)/2\in\mathbb Z\) y \(K_D\in\mathbb Z^3\). Fuera de ese dominio,
el tercer componente conserva el tipo \(\tfrac12\mathbb Z\).

\subsection{Naturalidad de los operadores bidireccionales y preservación de tipos}

El lector \(D_k\) es invariante bajo traslación temporal, reversión y toda
permutación biyectiva del alfabeto trítico. En las rutas declaradas, la
inversión simultánea de las dos orientaciones actúa por una traslación de
27 pasos discretos y conserva \(K_D\). El lector \(K_\Omega\) es invariante bajo
rotación y reversión dodecafásicas y bajo
\(C_M(\tau)=-\operatorname{rev}(\tau)\). Su forma cuadrática
\[
 \Omega(\tau)=\langle\tau,
 (4W_3^*W_3-3W_4^*W_4)\tau\rangle
\]
posee polarización bilineal entera y satisface la identidad de 2-cociclo.

Los contadores globales de giro, cambio de hoja, repetición de fase, cambio de
prefijo y retorno son constantes sobre las 324 rutas. Por consiguiente, todo
funcional que factorice exclusivamente por esos seis contadores es constante
y cancela en razones. Los lectores de
\eqref{eq:kappaD-general-rev11} y
\eqref{eq:kappaOmega-general-rev11} actúan sobre la historia local y superan
ese control de no factorización.

\begin{ablation}[Necesidad de las dos hojas en la ruta electrónica]
\label{abl:hojas-lector-two-way-rev11}
El lector cronológico produce
\[
 (A_D,B_D,C_D)=
 \begin{cases}
 (80,54,6),&\substack{\text{palabra cronológica conjunta de ambas hojas;}\\
 B_D=D_1(\Gamma_e)=54},\\
 (80,84,-8),&\text{proyección aditiva},\\
 (80,60,13),&\text{proyección multiplicativa}.
 \end{cases}
\]
La terna electrónica completa requiere la historia conjunta de las dos hojas
aritméticas.
\end{ablation}

\subsection{Censos exactos sobre las 324 rutas}

El lector cronológico produce catorce perfiles. La tabla siguiente registra su
censo completo; \(A_D=D_{27}+D_{36}\), \(B_D=D_1\) y
\(C_D=(D_4-D_3)/2\).

\begin{longtable}{@{}rrr r@{}}
\caption{Censo exacto de los catorce perfiles del lector cronológico sobre las
324 rutas orientadas.}
\label{tab:censo-perfiles-kd-rev11}\\
\toprule
\(A_D\)&\(B_D\)&\(C_D\)&multiplicidad\\
\midrule
\endfirsthead
\toprule
\(A_D\)&\(B_D\)&\(C_D\)&multiplicidad\\
\midrule
\endhead
0&24&0&18\\
40&24&2&36\\
40&54&6&18\\
40&78&27&36\\
40&84&30&36\\
60&78&25&18\\
60&84&28&18\\
80&54&6&18\\
80&54&7&18\\
80&66&2&18\\
80&66&3&36\\
80&66&4&18\\
100&66&0&18\\
100&66&2&18\\
\midrule
\multicolumn{3}{r}{total}&324\\
\bottomrule
\end{longtable}

El lector dodecafásico produce nueve perfiles:

\begin{table}[H]
\centering
\caption{Censo exacto de los nueve perfiles del lector dodecafásico sobre
las 324 rutas orientadas.}
\label{tab:censo-perfiles-komega-rev11}
\begin{tabular}{@{}rrr r@{}}
\toprule
\(\Omega\)&\(B_\Omega\)&\(P_6\)&multiplicidad\\
\midrule
80&54&6&198\\
56&54&5&68\\
48&54&5&34\\
36&54&4&8\\
20&54&4&4\\
4&54&2&4\\
40&54&4&4\\
24&54&4&2\\
24&54&2&2\\
\midrule
\multicolumn{3}{r}{total}&324\\
\bottomrule
\end{tabular}
\end{table}

Los lectores coinciden como triples en (18) rutas, siempre con el valor
\((80,54,6)\), y el par \(\mathbf K\) adopta (40) valores. En la sección
canónica \(B1\) de 81 posiciones, \(K_D\) es constante en 41 de las 56
familias y refina el cociente en 77 clases; \(K_\Omega\) es constante en 50
y produce 62 clases; el par produce 78. En el levantamiento de 324
orientaciones, \(K_D\) produce 146 clases y es no constante en todas las
familias; \(K_\Omega\) produce 100 y desciende en 17 familias; el par produce
151 y es no constante en todas. Estos censos prueban que las dos
proyecciones retienen memorias distintas y fijan el dominio de cada
afirmación de descenso.

\subsection{Trazas vectoriales y espectros conjuntos de las \(13\) multisecciones}

En la sección canónica \(B1\), denotamos por
\(\ell_j,\nu_j,d_j,u_j\) las multisecciones de generación \(j\) de los
sectores leptónico cargado, neutrino, quark de tipo abajo y quark de tipo
arriba, respectivamente. Para cada lector
\(r\in\{D,\Omega\}\) y cada componente \(a\in\{1,2,3\}\), definimos el
operador diagonal
\[
 K_{M,a}^r e_\Gamma=K_{r,a}(\Gamma)e_\Gamma,
 \qquad
 \mathbf K_M^r=(K_{M,1}^r,K_{M,2}^r,K_{M,3}^r).
\]
Los tres operadores de \(\mathbf K_M^r\) son autoadjuntos y conmutan. Por
tanto, su traza normalizada se entiende componente a componente,
\[
 \overline{\mathbf K}_M^r
 =\frac{1}{\dim M}
 \bigl(\operatorname{Tr}K_{M,1}^r,
       \operatorname{Tr}K_{M,2}^r,
       \operatorname{Tr}K_{M,3}^r\bigr),
\]
y su espectro es el espectro conjunto, es decir, la distribución con
multiplicidad de los triples simultáneos. Las trece trazas vectoriales son:

\begin{longtable}{@{}c r c c@{}}
\caption{Trazas vectoriales normalizadas de los lectores \(K_D\) y
\(K_\Omega\) en las trece multisecciones canónicas.}
\label{tab:trazas-multisecciones-two-way-rev11}\\
\toprule
\(M\)&\(\dim M\)&\(\overline{\mathbf K}_M^D\)&\(\overline{\mathbf K}_M^\Omega\)\\
\midrule
\endfirsthead
\toprule
\(M\)&\(\dim M\)&\(\overline{\mathbf K}_M^D\)&\(\overline{\mathbf K}_M^\Omega\)\\
\midrule
\endhead
\(\ell_0\)&1&\((80,54,6)\)&\((80,54,6)\)\\
\(\ell_2\)&8&\((65,249/4,17/2)\)&\((80,54,6)\)\\
\(\ell_4\)&16&\((225/4,489/8,21/2)\)&\((141/2,54,11/2)\)\\
\(\nu_1\)&2&\((40,51,29/2)\)&\((60,54,5)\)\\
\(\nu_2\)&4&\((45,63,29/2)\)&\((61,54,5)\)\\
\(\nu_3\)&6&\((170/3,55,34/3)\)&\((206/3,54,11/2)\)\\
\(\nu_4\)&8&\((105/2,237/4,95/8)\)&\((143/2,54,45/8)\)\\
\(d_1\)&2&\((40,51,29/2)\)&\((60,54,5)\)\\
\(d_2\)&4&\((45,63,57/4)\)&\((61,54,5)\)\\
\(d_3\)&6&\((170/3,55,23/2)\)&\((218/3,54,17/3)\)\\
\(d_4\)&8&\((105/2,237/4,49/4)\)&\((149/2,54,23/4)\)\\
\(u_1\)&4&\((65,66,9)\)&\((80,54,6)\)\\
\(u_3\)&12&\((205/3,119/2,113/12)\)&\((74,54,23/4)\)\\
\bottomrule
\end{longtable}

La notación \((a,b,c)^{\times n}\) indica multiplicidad \(n\). El espectro
conjunto completo de \(\mathbf K_D\) es:

\begingroup\footnotesize
\begin{longtable}{@{}c p{.82\textwidth}@{}}
\caption{Espectros conjuntos completos del lector \(K_D\) en las trece
multisecciones canónicas.}
\\
\toprule
\(M\)&\(\operatorname{Spec}_{\rm conj}\mathbf K_M^D\)\\
\midrule
\endhead
\(\ell_0\)&\((80,54,6)^{\times1}\)\\
\(\ell_2\)&\((0,24,0)^{\times1};(40,78,27)^{\times2};(80,54,7)^{\times1};(80,66,2)^{\times1};(80,66,3)^{\times1};(100,66,0)^{\times1};(100,66,2)^{\times1}\)\\
\(\ell_4\)&\((0,24,0)^{\times1}\); \((40,24,2)^{\times2}\); \((40,54,6)^{\times2}\); \((40,84,30)^{\times2}\); \((60,78,25)^{\times3}\); \((80,66,2)^{\times2}\); \((80,66,3)^{\times3}\); \((80,66,4)^{\times1}\)\\
\(\nu_1\)&\((40,24,2)^{\times1};(40,78,27)^{\times1}\)\\
\(\nu_2\)&\((0,24,0)^{\times1};(40,84,30)^{\times1};(60,78,25)^{\times1};(80,66,3)^{\times1}\)\\
\(\nu_3\)&\((40,24,2)^{\times2};(40,78,27)^{\times1};(60,84,28)^{\times1};(80,54,6)^{\times1};(80,66,3)^{\times1}\)\\
\(\nu_4\)&\((0,24,0)^{\times1};(40,24,2)^{\times1};(40,78,27)^{\times2};(40,84,30)^{\times1};(80,54,7)^{\times1};(80,66,2)^{\times1};(100,66,0)^{\times1}\)\\
\(d_1\)&\((40,24,2)^{\times1};(40,78,27)^{\times1}\)\\
\(d_2\)&\((0,24,0)^{\times1};(40,84,30)^{\times1};(60,78,25)^{\times1};(80,66,2)^{\times1}\)\\
\(d_3\)&\((40,24,2)^{\times2};(40,78,27)^{\times1};(60,84,28)^{\times1};(80,54,6)^{\times1};(80,66,4)^{\times1}\)\\
\(d_4\)&\((0,24,0)^{\times1};(40,24,2)^{\times1};(40,78,27)^{\times2};(40,84,30)^{\times1};(80,54,7)^{\times1};(80,66,3)^{\times1};(100,66,2)^{\times1}\)\\
\(u_1\)&\((40,54,6)^{\times1};(60,78,25)^{\times1};(80,66,2)^{\times1};(80,66,3)^{\times1}\)\\
\(u_3\)&\((40,24,2)^{\times2}\); \((40,78,27)^{\times2}\); \((60,84,28)^{\times1}\); \((80,54,6)^{\times3}\); \((80,66,3)^{\times1}\); \((80,66,4)^{\times1}\); \((100,66,0)^{\times1}\); \((100,66,2)^{\times1}\)\\
\bottomrule
\end{longtable}
\label{tab:espectros-kd-multisecciones-rev11}
\endgroup

El espectro conjunto completo de \(\mathbf K_\Omega\) sobre las mismas fibras es:

\begingroup\footnotesize
\begin{longtable}{@{}c p{.82\textwidth}@{}}
\caption{Espectros conjuntos completos del lector \(K_\Omega\) en las trece
multisecciones canónicas.}
\label{tab:espectros-komega-multisecciones-rev11}\\
\toprule
\(M\)&\(\operatorname{Spec}_{\rm conj}\mathbf K_M^\Omega\)\\
\midrule
\endfirsthead
\toprule
\(M\)&\(\operatorname{Spec}_{\rm conj}\mathbf K_M^\Omega\)\\
\midrule
\endhead
\(\ell_0\)&\((80,54,6)^{\times1}\)\\
\(\ell_2\)&\((80,54,6)^{\times8}\)\\
\(\ell_4\)&\((24,54,2)^{\times1};(56,54,5)^{\times4};(80,54,6)^{\times11}\)\\
\(\nu_1\)&\((40,54,4)^{\times1};(80,54,6)^{\times1}\)\\
\(\nu_2\)&\((4,54,2)^{\times1};(80,54,6)^{\times3}\)\\
\(\nu_3\)&\((36,54,4)^{\times1};(56,54,5)^{\times1};(80,54,6)^{\times4}\)\\
\(\nu_4\)&\((36,54,4)^{\times1};(56,54,5)^{\times1};(80,54,6)^{\times6}\)\\
\(d_1\)&\((40,54,4)^{\times1};(80,54,6)^{\times1}\)\\
\(d_2\)&\((4,54,2)^{\times1};(80,54,6)^{\times3}\)\\
\(d_3\)&\((36,54,4)^{\times1};(80,54,6)^{\times5}\)\\
\(d_4\)&\((36,54,4)^{\times1};(80,54,6)^{\times7}\)\\
\(u_1\)&\((80,54,6)^{\times4}\)\\
\(u_3\)&\((56,54,5)^{\times3};(80,54,6)^{\times9}\)\\
\bottomrule
\end{longtable}
\endgroup

Las trazas pueden coincidir entre multisecciones de tipo diferente; el tipo,
la multiplicidad y el espectro completo permanecen como datos del operador.
El certificado reproduce estas tablas fila por fila.

\HMTClaim{MD-R-MASS-FIBRE-OPERATORS}
\subsection{Operadores de fibra y naturalidad del cambio de base}

Distinguimos la fibra canónica \(F_M^{B1}\), formada por las posiciones de
la multisección \(M\) con su orientación publicada, y la fibra orientada
\(F_M^{\rm or}\), formada por sus cuatro levantamientos por posición. Para
\(s\in\{B1,{\rm or}\}\), sea
\(\mathcal H_{M,s}=\bigoplus_{\Gamma\in F_M^s}\mathbb C e_\Gamma\). Definimos
\begin{equation}
 B_{M,s}^r e_\Gamma=\kappa_r(\Gamma)e_\Gamma,
 \qquad
 R_{\rm act}^{B_{M,s}^r}
 =\exp\bigl[(\log R_{\rm act})B_{M,s}^r\bigr],
 \qquad r\in\{D,\Omega\}.
 \label{eq:operadores-fibra-two-way-rev11}
\end{equation}
Cada \(B_{M,s}^r\) es autoadjunto y diagonal en la base de rutas, por lo que
\(R_{\rm act}^{B_{M,s}^r}\) es positivo. Para una masa basal positiva
\(\widehat M_{M,s}^{(0)}\), la forma operatoria condicionada a cada lector es
\begin{equation}
 \widehat M_{M,s}^{\beta,r}
 =R_{\rm act}^{B_{M,s}^r/2}\,
  \widehat M_{M,s}^{(0)}\,
  R_{\rm act}^{B_{M,s}^r/2}.
 \label{eq:ley-operatoria-fibra-rev11}
\end{equation}
Esta expresión es positiva sin imponer conmutatividad. Cuando
\([\widehat M_{M,s}^{(0)},B_{M,s}^r]=0\), reduce a
\(\widehat M_{M,s}^{(0)}R_{\rm act}^{B_{M,s}^r}\).
Un cambio de base unitario \(U\) transforma
\(B_{M,s}^r\mapsto UB_{M,s}^rU^*\) y conserva espectro, traza, determinante y
polinomio característico. La fibra electrónica canónica \(F_e^{B1}\) tiene
rango uno y reduce al escalar del
\cref{thm:reduccion-electron-two-way-rev11}.
Su levantamiento orientado tiene rango cuatro: las rutas \((++)\) y \((--)\)
comparten el perfil electrónico, mientras las orientaciones mixtas conservan
perfiles distintos.

El determinante normalizado
\[
 \operatorname{ndet}(R_{\rm act}^{B_{M,s}^r})
 =R_{\rm act}^{\operatorname{Tr}(B_{M,s}^r)/\dim F_M^s}
\]
es natural bajo permutación de la fibra. Aditividad logarítmica lo distingue
como media geométrica, pero no se usa como sustituto de la realización física.
La salida exacta de refinamiento sigue siendo el par
\((B_{M,s}^D,B_{M,s}^\Omega)\); la sección escalar de una especie se obtiene
por el morfismo de sabor--generación--ruta ya construido en
\cref{eq:morfismo-sabor-ruta-cerrado-rev2,eq:realizacion-fisica-sabor-ruta-cerrada-rev2}:
selecciona la ruta persistente por datos no metrológicos y aplica después el
carácter de su firma. Una funcional escalar arbitraria
\(S(B_{M,s}^D,B_{M,s}^\Omega)\) no define esa selección y permanece sólo como
control de que no se reintroduce la masa objetivo.

\subsection{Escala absoluta, corrección relativa y sensibilidades}

La genealogía de la recta de acción se escribe
\[
 \Sigma_H=\varphi\alpha_{\rm HMT}^{16},
 \qquad \operatorname{ord}_{10}(\Sigma_H)=-34,
\]
\[
 \delta_{2w}=H_5-\eta_{\rm ret}
 =\frac{90}{\pi}\mathscr C_\pi(\alpha_{\rm HMT}),
 \qquad R_{\rm act}=\frac{H_5}{\eta_{\rm ret}},
\]
\[
 \hbar_{\rm pre}=H_5\,10^{-34}\mathcal U_S,
 \qquad
 \hbar_{\rm ret}=\eta_{\rm ret}\,10^{-34}\mathcal U_S.
\]
El factor \(10^{-34}\mathcal U_S\) fija la sección absoluta y cancela en
razones. Para el lector dodecafásico,
\begin{equation}
 \log\!\left[
 \frac{(m_X^{\beta,\Omega}/m_e^{\beta})}
 {(m_X^{(0)}/m_e^{(0)})}\right]
 =\bigl[(\Omega_X-80)+(P_{6,X}-6)\Delta_4\bigr]
 \log R_{\rm act}.
 \label{eq:razon-masa-omega-rev11}
\end{equation}
El término común \(-54\alpha_{\rm HMT}\) cancela también. Para el lector
cronológico,
\begin{equation}
 \log\!\left[
 \frac{(m_X^{\beta,D}/m_e^{\beta})}
 {(m_X^{(0)}/m_e^{(0)})}\right]
 =\bigl[\kappa_D(\Gamma_X)-\kappa_e\bigr]\log R_{\rm act}.
 \label{eq:razon-masa-D-rev11}
\end{equation}
Con \(\log R_{\rm act}=6.932817628081925\ldots\times10^{-10}\),
\[
 \frac{\partial\log m}{\partial\Omega}=6.9328\times10^{-10},
 \quad
 \frac{\partial\log m}{\partial P_6}
 =\Delta_4\log R_{\rm act}=7.7090\times10^{-14},
\]
\[
 \frac{\partial\log m}{\partial\alpha}
 =-54\log R_{\rm act}=-3.74372\times10^{-8},
 \quad
 \left.\frac{\partial\log m}{\partial\Delta_4}\right|_e
 =6\log R_{\rm act}=4.15969\times10^{-9}.
\]
Estas derivadas son sensibilidades del modelo y permanecen separadas de la
incertidumbre experimental de una masa de contraste.

\subsection{Emparejamiento físico y sector nulo}

El emparejamiento tipado sigue la cadena
\[
 \begin{aligned}
 &\text{registro físico tipado}
 \longrightarrow\text{multisección o representación}
 \longrightarrow\text{fibra de rutas}\longrightarrow\text{registro}\\
 &\longrightarrow\text{firma y carácter}
 \longrightarrow(B_{M,s}^D,B_{M,s}^\Omega)
 \longrightarrow\text{carta dimensional}
 \longrightarrow\text{contraste externo}.
 \end{aligned}
\]
El inventario físico tipado distingue registros por objeto, representación y codominio; las 324 rutas son
el motor interno; las 471 filas del catálogo PDG son propiedades externas.
Estos cardinales pertenecen a dominios diferentes. La flecha se construye
desde descriptores físicos no metrológicos, la multisección y la ruta; una
búsqueda por proximidad a una masa pertenece al régimen de calibración.

El sector nulo se bifurca antes del carácter positivo. Para fotón y gluones,
\[
 \kappa_{\rm null}=\mathrm{N/A},
 \qquad m_{\rm null}=0.
\]
La expresión \(R_{\rm act}^{0}=1\) es la unidad del carácter positivo. Los
compuestos concatenan
registros antes de evaluar la firma; los sectores topológicos conservan fusión y
trenza; una sección virtual se clasifica por su horizonte finito de
prolongación.

\subsection{Independencia del constructor de rutas respecto de las masas y clausura del operador}

La verificación estructural utiliza exclusivamente el censo de rutas, el
registro \(81\times108\) y el transductor de las 81 posiciones. Comprueba su
acoplamiento fila a fila, la cobertura
\(81/56/13/324\), las naturalidades, el triple electrónico \((80,54,6)\), los
censos (14) y (9), la coincidencia (18/324), los
(40) pares, descensos diferenciados en \(B1\) y en las 324 orientaciones,
periodicidad 54, inversión por desplazamiento 27 y ablaciones de hoja. El
sector nulo pertenece al dominio exterior al carácter positivo. El cálculo se
realiza antes de introducir masas, CODATA, PDG, residuos metrológicos o
etiquetas nominales.

Este certificado no evalúa \(\alpha_{\rm HMT}\), \(\Delta_4\),
\(R_{\rm act}\), \(\kappa_e\) ni la carta MeV. Esas evaluaciones pertenecen a
la cadena aritmética de la recta de acción y a la realización energética del
\cref{thm:reduccion-electron-two-way-rev11}; esta cadena aritmética permanece
separada de la comprobación combinatoria de los lectores.

La comparación metrológica aplica después los operadores al inventario tipado,
declara unidad, esquema, escala, fecha e incertidumbre, y ejecuta permutaciones
y validación por retención de multisecciones completas. La salida física queda
gobernada por el dominio \(81/56/13/324\), por los dos lectores de ruta y por
los operadores de fibra.

\begin{statusbox}[title={Dominio exacto y realización física cerrada}]
El atlas, la cadena \(C_{81}\to\mathcal L_{108}\) del registro cronológico
enriquecido de \(108\) iteraciones,
la recta de acción,
los lectores \(K_D,K_\Omega\), su reducción electrónica y los operadores de
fibra son resultados internos exactos. El par
\(\mathbf K=(K_D,K_\Omega)\) es la salida completa sobre cada ruta y
determina los dos operadores diagonales de toda fibra. En la fibra
electrónica ambos lectores coinciden. En las fibras no centrales, el par y su
espectro preservan el refinamiento completo, mientras
\(\mathcal S_{\rm phys}\) liga rama, generación, firma fina y, cuando
corresponde, órbita de color con la ruta persistente. El carácter de esa ruta
publica la coordenada escalar antes del contraste metrológico. Por ello no hay
una disyunción entre operador cerrado y masa pendiente: son dos lecturas
tipadas de la misma genealogía. Ambos lectores son co-rectores, y ninguna
proximidad a una masa externa interviene en la selección.
\end{statusbox}
\HMTClaim{MD-R-MASS-PHYSICAL-SELECTION-CLOSED}
